\chapter{Proposed Pipeline: SMVMD Common Modes, Riemannian Geometry and Gradient Boosting}
\thispagestyle{empty}
\label{ch:pipeline}

\section{Overview}
The proposed pipeline processes every 5-second window $\mathbf{X}\in\mathbb{R}^{C\times T}$
($C$ channels, $T=640$ samples) in five stages (Fig.~\ref{fig:pipeline}):
\begin{enumerate}
\item \textbf{SMVMD} splits the window into a window-specific number $K$ of multivariate modes,
      each concentrated around one centre frequency shared by all channels;
\item \textbf{common modes}: the modes are grouped by centre frequency into four fixed bands
      ($\delta$, $\theta$, $\alpha$, $\beta$) and summed within each band, so that every window
      is described by the same four components;
\item \textbf{covariance}: each common mode is summarised by its $C\times C$ channel covariance
      matrix;
\item \textbf{Riemannian tangent space}: the covariance matrices are mapped to vectors in the
      tangent space at their Riemannian mean, computed on the training windows;
\item \textbf{gradient-boosted trees} (LightGBM, or XGBoost) classify the concatenated tangent
      vectors; the windows of a person are combined by majority vote.
\end{enumerate}
Stages 1--3 operate on one window at a time and use no labels and no other windows, so they
cannot leak information between people. Everything that is learned from data (the reference
point of the tangent space, the standardisation, the classifier, its depth and the set of modes
used) is fitted on the training subjects of each fold only.

\begin{figure}[H]\centering
\begin{tikzpicture}[node distance=4mm and 4mm, font=\footnotesize,
  box/.style={draw, rounded corners=2pt, align=center, minimum height=10mm, text width=21mm,
              fill=gray!6},
  arr/.style={-{Stealth[length=2mm]}, thick}]
\node[box] (x) {EEG window\\$C\times 640$};
\node[box, right=of x] (s) {SMVMD\\$K$ modes,\\centre freq.\ $f_k$};
\node[box, right=of s] (c) {common modes\\$\delta,\theta,\alpha,\beta$\\by $f_k$};
\node[box, right=of c] (v) {covariance\\per mode\\(OAS)};
\node[box, below=9mm of v] (t) {Riemannian\\tangent space\\(train mean)};
\node[box, left=of t] (g) {LightGBM /\\XGBoost};
\node[box, left=of g] (m) {majority vote\\per person};
\node[box, left=of m, fill=white] (d) {IDD / ADHD\\or control};
\draw[arr] (x) -- (s); \draw[arr] (s) -- (c); \draw[arr] (c) -- (v); \draw[arr] (v) -- (t);
\draw[arr] (t) -- (g); \draw[arr] (g) -- (m); \draw[arr] (m) -- (d);
\end{tikzpicture}
\caption{The proposed pipeline. The upper row is computed independently for each window and uses
no labels; the lower row is learned from the training subjects of each fold.}
\label{fig:pipeline}
\end{figure}

\section{Successive multivariate variational mode decomposition}
\textbf{Variational mode decomposition.} VMD~\cite{dragomiretskiy2014} represents a signal $x(t)$
as a sum of $K$ modes $u_k(t)$, each compact around a centre frequency $\omega_k$, by solving
\begin{equation}
\min_{\{u_k\},\{\omega_k\}} \sum_{k=1}^{K}
\left\| \partial_t\!\left[\left(\delta(t)+\frac{j}{\pi t}\right) * u_k(t)\right] e^{-j\omega_k t}
\right\|_2^2
\quad\text{s.t.}\quad \sum_{k} u_k(t) = x(t).
\end{equation}
The term in brackets is the analytic signal of the mode; shifting it by $e^{-j\omega_k t}$ to
baseband and penalising the energy of its derivative measures the bandwidth of the mode. In the
Fourier domain the solution for a mode is a Wiener-type filter of the signal centred at
$\omega_k$, and the centre frequency is the power-weighted mean frequency of the mode.

\textbf{Multivariate extension.} For a $C$-channel signal, MVMD~\cite{rehman2019} looks for
multivariate modes $\mathbf{u}_k(t)\in\mathbb{R}^C$ whose channels share one centre frequency
$\omega_k$. This alignment is essential here: the $c$-th channel of mode $k$ describes the same
oscillation at every electrode, so the covariance between channels within one mode is meaningful.

\textbf{Successive extraction.} VMD and MVMD need $K$ in advance. SMVMD~\cite{liu2022} extracts
the modes one at a time from the residual: the first mode is extracted from the window, then the
next from what remains, and so on, until the residual energy falls below 1\% of the window
energy or 20 modes have been extracted. In the implementation used here (the same as for the
starting method), mode $k$ is obtained from the residual spectrum $\hat{r}_c(\nu)$ of each
channel $c$, with $\nu\in[0,\tfrac12]$ the normalised frequency in cycles per sample, by
iterating
\begin{equation}
\hat{u}_{k,c}(\nu) = \frac{\hat{r}_c(\nu)}{1 + 2\alpha(\nu-\nu_k)^2},
\qquad
\nu_k = \frac{\sum_c \sum_{\nu} \nu\, |\hat{u}_{k,c}(\nu)|^2}
             {\sum_c \sum_{\nu} |\hat{u}_{k,c}(\nu)|^2},
\label{eq:smvmd}
\end{equation}
for at most 500 iterations or until $\nu_k$ changes by less than $10^{-10}$, after which the mode
is subtracted from the residual. The common centre frequency $\nu_k$, shared by all channels, is
what makes the modes multivariate. The parameter $\alpha$ sets the bandwidth of the modes ($\alpha=2000$ for DS-1 and
$1000$ for DS-2, as for the starting method~\cite{chandela2025}), and the dual-ascent step is set
to zero, so the modes do not have to reconstruct the noisy signal exactly. Because the number of
modes adapts to the signal, it differs between windows: on average \R{CmDSoneRestMusicModesMean}
modes per DS-1 window (range \R{CmDSoneRestMusicModesMin}--\R{CmDSoneRestMusicModesMax}) and
\R{CmDStwoModesMean} per DS-2 window (range \R{CmDStwoModesMin}--\R{CmDStwoModesMax}).

\section{Common modes}
A classifier needs the same features for every window, but SMVMD gives each window its own
number of modes at its own frequencies (Fig.~\ref{fig:centrefreq}). The starting method solved
this by averaging the features of all modes with energy weights, which discards the frequency
information. The proposed pipeline instead assigns every mode to a fixed frequency band by its
centre frequency in hertz, $f_k=\nu_k f_s$ with $f_s=128$~Hz, and sums the modes within each
band:
\begin{equation}
\mathbf{U}_b = \sum_{k:\; f_k\in B_b} \mathbf{u}_k,\qquad
B_\delta=[0,4),\; B_\theta=[4,8),\; B_\alpha=[8,13),\; B_\beta=[13,30)\ \text{Hz}.
\label{eq:common}
\end{equation}
The four components $\mathbf{U}_\delta,\mathbf{U}_\theta,\mathbf{U}_\alpha,\mathbf{U}_\beta \in
\mathbb{R}^{C\times T}$ are the \emph{common modes} of the window. Unlike the outputs of fixed
band-pass filters, they are built from the window's own oscillatory components: a common mode
contains the modes that SMVMD found in that band, with the bandwidth and centre that fit the
window, and nothing where SMVMD found no oscillation. The band edges are the conventional EEG
bands, chosen before any result was seen. Modes at 30~Hz and above are not used: DS-1 was
band-limited to 1--30~Hz by its authors, and in DS-2 such modes carry on average only
\R{CmDStwoShareAbove}\% of the mode energy. A fifth component, the sum of all modes (the SMVMD
reconstruction), is used only in the ablation.

\begin{figure}[H]\centering
\genfig[0.95\textwidth]{centre_frequencies}
\caption{Centre frequencies of all SMVMD modes of all windows. Dashed lines: the edges of the
common modes (4, 8, 13 and 30~Hz).}
\label{fig:centrefreq}
\end{figure}

Table~\ref{tab:commonmodes} summarises the common modes. In DS-1 the energy is shared between the
$\delta$ and $\alpha$ modes, consistent with relaxed wakefulness during the rest and music
recordings. In DS-2 nearly all energy lies in the $\delta$ mode, and many windows contain no mode
in the $\alpha$ (\R{CmDStwoEmptyAlpha}\% of windows) or $\beta$ (\R{CmDStwoEmptyBeta}\%) band.

\begin{table}[H]
\caption{Common modes: modes per window (mean and range), and per band the mean share of mode
energy (\%) / the percentage of windows without any mode in that band.}
\label{tab:commonmodes}\centering
\gentable{tab_common_modes}
\end{table}

\section{Covariance estimation}
Each common mode is summarised by its spatial covariance matrix. With
$\mathbf{S}_b=\frac{1}{T}\mathbf{U}_b\mathbf{U}_b^{\top}$ the sample covariance (the modes have
zero mean), the oracle approximating shrinkage (OAS) estimator~\cite{chen2010}
\begin{equation}
\boldsymbol{\Sigma}_b = (1-\rho)\,\mathbf{S}_b + \rho\,\frac{\operatorname{tr}(\mathbf{S}_b)}{C}\,\mathbf{I}
\end{equation}
shrinks it towards a scaled identity, with the weight $\rho$ computed in closed form from
$\mathbf{S}_b$. Shrinkage keeps the matrix well conditioned, which matters for DS-1, whose
cleaned recordings are rank-deficient after the removal of independent components. When a window
has no mode in a band, $\mathbf{U}_b=\mathbf{0}$; to keep every matrix positive definite, a floor
$\epsilon\,\bar{p}\,\mathbf{I}$ is added to all matrices, where $\bar{p}$ is the mean channel power
of the whole window and $\epsilon=10^{-4}$. The covariance of a window depends only on that
window, so these matrices are computed once, before cross-validation.

The covariance matrix captures two kinds of information: the diagonal holds the power of each
channel in the band, and the off-diagonal entries hold the co-variation of pairs of channels,
that is, the spatial structure of the oscillation. Per-channel features such as band powers
contain only the first kind.

\section{Riemannian geometry of covariance matrices}
\textbf{The SPD manifold.} Covariance matrices are symmetric positive definite (SPD). They do not
form a vector space: the difference of two SPD matrices need not be SPD, and the Euclidean mean
of SPD matrices is inflated (the ``swelling effect''). The set of SPD matrices is a Riemannian
manifold, and the natural distance on it, the affine-invariant Riemannian distance,
\begin{equation}
\delta_R(\boldsymbol{\Sigma}_1,\boldsymbol{\Sigma}_2) =
\left\| \log\!\left(\boldsymbol{\Sigma}_1^{-1/2}\boldsymbol{\Sigma}_2\boldsymbol{\Sigma}_1^{-1/2}\right)\right\|_F
= \Big(\sum_{i=1}^{C} \log^2 \lambda_i\Big)^{1/2},
\end{equation}
where $\lambda_i$ are the eigenvalues of $\boldsymbol{\Sigma}_1^{-1}\boldsymbol{\Sigma}_2$, is
invariant to any invertible linear transformation of the channels:
$\delta_R(\mathbf{A}\boldsymbol{\Sigma}_1\mathbf{A}^\top, \mathbf{A}\boldsymbol{\Sigma}_2\mathbf{A}^\top)
=\delta_R(\boldsymbol{\Sigma}_1,\boldsymbol{\Sigma}_2)$~\cite{barachant2012, congedo2017}. Since
the scalp EEG is a linear mixture of cortical sources, this invariance makes the geometry well
suited to EEG.

\textbf{Riemannian mean.} The mean of a set of matrices $\{\boldsymbol{\Sigma}_i\}$ is the point
$\mathbf{M}$ that minimises $\sum_i \delta_R^2(\mathbf{M},\boldsymbol{\Sigma}_i)$. It has no closed
form and is computed by a fixed-point iteration that alternately maps the matrices to the tangent
space at the current estimate, averages them there and maps the average back to the manifold.

\textbf{Tangent-space mapping.} Near a reference point $\mathbf{M}$, the manifold is approximated
by its tangent space, which is a vector space. Each matrix is mapped to a vector by
\begin{equation}
\mathbf{s} = \operatorname{upper}\!\left(\log\!\left(\mathbf{M}^{-1/2}\boldsymbol{\Sigma}\,\mathbf{M}^{-1/2}\right)\right),
\label{eq:ts}
\end{equation}
where $\operatorname{upper}(\cdot)$ keeps the $C(C+1)/2$ entries of the upper triangle and
multiplies the off-diagonal entries by $\sqrt{2}$, so that the Euclidean norm of $\mathbf{s}$
equals the Riemannian distance between $\boldsymbol{\Sigma}$ and $\mathbf{M}$. The reference
$\mathbf{M}$ is the Riemannian mean of the \emph{training} windows of the fold, computed
separately for each common mode. The four tangent vectors are concatenated into one feature
vector of length $4\cdot C(C+1)/2$: \R{RtFeatures} features for DS-2 ($C=19$) and 420 for DS-1
($C=14$).

\section{Gradient-boosted decision trees}
The tangent vectors are classified with gradient-boosted decision trees~\cite{chen2016, ke2017}.
The model is an additive ensemble $F(\mathbf{s})=\sum_{m=1}^{M} \eta\, f_m(\mathbf{s})$ of $M$
regression trees, and the probability of the disorder is $\sigma(F(\mathbf{s}))$. Trees are
added one at a time; tree $f_m$ is fitted to the gradient (and, in XGBoost and LightGBM, the
curvature) of the logistic loss of the current ensemble, so that each tree corrects the errors of
the previous ones. Gradient boosting was chosen over the logistic regression and SVM used in the
preliminary study for three reasons: trees model interactions between features, here between
channel pairs and between modes, without a kernel choice; they are insensitive to the scale and
to outlying values of individual features, such as the floored covariances of empty modes; and
they train in seconds and give a feature importance that can be traced back to modes and
electrodes.

\textbf{LightGBM}~\cite{ke2017} grows trees leaf-wise on histogram-binned features, and
\textbf{XGBoost}~\cite{chen2016} grows them level-wise with a regularised objective. Both are
used with the same budget, fixed in advance: $M=300$ trees, learning rate $\eta=0.05$, a random
80\% of the windows and 50\% of the features per tree, and class weights inversely proportional
to the class frequencies of the training windows. The tree depth is chosen from $\{2, 4\}$ by the
inner cross-validation (for LightGBM the number of leaves is $2^{\text{depth}}-1$). The features
are standardised with the mean and standard deviation of the training windows.

\section{Choice of the common modes}
Which common modes should be used may differ between disorders, so the mode set is part of the
model and is chosen by the inner cross-validation, together with the tree depth. The candidates
are each single mode and all four together, giving $5\times 2=10$ grid points. In each outer
fold, an inner stratified grouped 4-fold cross-validation over the training subjects evaluates
every grid point, the point with the highest pooled balanced accuracy is selected, and the model
is refitted on the whole outer training set. The test subjects of the outer fold play no part in
this choice.

\section{Decision per person}
The pipeline classifies every window. A person is assigned the class of the majority of their
windows, with ties broken by the mean predicted probability.

\section{What is fitted where}
\begin{table}[H]
\caption{Where each part of the pipeline is computed. ``Per window'' parts use one window and no
labels; ``training subjects'' parts are refitted in every outer fold.}
\label{tab:fitted}\centering\small
\begin{tabular}{lll}
\hline
Stage & Computed from & Leakage possible? \\
\hline
SMVMD, common modes & each window alone & no \\
OAS covariance and floor & each window alone & no \\
Riemannian mean (tangent-space reference) & training subjects of the fold & no \\
Standardisation & training subjects of the fold & no \\
Mode set and tree depth & inner CV on training subjects & no \\
Gradient-boosted trees & training subjects of the fold & no \\
Majority vote & the test person's own predictions & no \\
\hline
\end{tabular}
\end{table}

\section{Computational cost}
\label{sec:cost}
On the laptop CPU used for this work (one core per stage), processing one 5-second window takes
about \R{RtSmvmd}~ms for SMVMD, \R{RtCov}~ms for the four covariance matrices and \R{RtPred}~ms for
the tangent-space mapping and LightGBM prediction, \R{RtTotal}~ms in total, about
\R{RtRealtime} times faster than real time. Training one model on \R{RtTrainWin} DS-2 windows
takes about \R{RtTrain}~s; the trained model has \R{RtTrees} trees with \R{RtLeaves} leaves in
total and occupies about \R{RtModelKB}~kB. No GPU is needed at any stage. SMVMD dominates the
cost; in the experiments it was run in parallel over windows.

\section{Implementation}
The pipeline is implemented in Python with NumPy and SciPy (SMVMD, filters), pyRiemann (OAS
covariance, Riemannian mean, tangent space), scikit-learn~\cite{pedregosa2011} (standardisation,
logistic regression, SVM, cross-validation), XGBoost and LightGBM. The experiments are in
\texttt{experiments/stage2c\_smvmd\_riemann.py} and \texttt{experiments/stage2d\_main\_pipeline.py}
(Appendix~B).
